%% ===================================================================
\section{\vrev{학술적 기여}}
\label{sec:conclusion-contributions}
%% ===================================================================


%% -------------------------------------------------------------------
\subsection{\vrev{이론적 기여}}
\label{subsec:contrib-theoretical}
%% -------------------------------------------------------------------

\fdel{본 연구의 이론적 기여는 다음의 세 가지로 정리된다.}

\fdel{첫째, \textbf{업무연속성관리$\cdot$관측성$\cdot$시그모이드 바운딩의 통합적 이론 기반을
확립}하였다.
기존 BCM 연구에서 RTO는 BIA 시점에 고정되는 정적 매개변수로 취급되었으며,
관측성과 불확실성을 단일 수리 모델에서 통합한 선행 연구는 부재하였다.
본 연구는 제어 이론의 관측성 개념과 확률론적 불확실성 모형을
수정 시그모이드 함수라는 공통 기저 함수 위에서 결합함으로써,
BCM 분야에 새로운 이론적 기반을 제공한다.
특히 $\DRTO = R_0 + E_{O,\text{bnd}} + E_{U,\text{bnd}}$의 가산 구조는
각 효과의 독립적 분석과 결합 효과의 동시 포착을 가능하게 하여,
모델의 해석 가능성(interpretability)과 확장 가능성(extensibility)을
동시에 확보한다.}

\fdel{둘째, \textbf{꼬리 감쇠 현상을 발견하고 정량화}하였다.
P85.8 분할 기준을 경계로 관측성 가중치 $k_O$의 효과가
Core($-0.798$)에서 Tail($-0.415$)로 $0.52$배 감쇠되는 현상은,
극단 불확실성 환경에서 시그모이드 포화에 의해
관측성 채널의 조절 효과가 구조적으로 제한됨을 규명한 것이다.
이 발견은 BCM 자원 배분 전략에 새로운 분석적 기반을 제시하며,
극단 상황에서의 대비 전략이 관측성 투자만으로는 한계가 있으며
사전 자원 확보가 병행되어야 한다는 실무적 시사점을 제공한다.}

\fdel{셋째, \textbf{P85.8 CDF 교차점을 가우시안--파레토 분포 전환의 구조적 경계로 규명}하였다.
C2와 C3의 CDF가 P85.8($\eta \approx 2.42$)에서 교차하며,
이 교차점을 경계로 두 조건의 행태가 질적으로 전환된다.
P85.8 이하에서 두 분포는 사실상 구분 불가능하나,
P85.8 초과에서 C3의 $\eta$가 C2를 급격히 상회하여
CVaR$_{99}$ 기준 $+24.0\%$의 격차를 보인다.
이는 파레토 불확실성의 효과가 \textit{평균 수준의 비차별성}과
\textit{꼬리 영역의 극적 분기}로 구성되는 이중 구조를 가진다는 것을
실증적으로 규명한 결과이다.}

